\documentclass[conference]{IEEEtran}
\usepackage{cite}
\usepackage{amsmath,amssymb,amsfonts}
\usepackage{array}
\usepackage{booktabs}
\usepackage{url}

\begin{document}

\title{Embedded Radio Frequency and Computer Vision-Based Spy Camera Detector with Smartphone Integration}

\author{\IEEEauthorblockN{Samarth S. Prabhu, Saamarth S., Amogh Neeli, Prajwal V. Jois, Dr. C. Nandeesh Kumar, and Dr. Pratiba D.}
\IEEEauthorblockA{Rashtreeya Vidyalaya College of Engineering, Bengaluru, India}}

\maketitle

\begin{abstract}
[Abstract reserved for Milestone 4.]
\end{abstract}

\begin{IEEEkeywords}
[Index terms reserved for Milestone 4.]
\end{IEEEkeywords}

\section{Introduction}
Miniaturized cameras can be concealed inside ordinary room fixtures, including clocks, chargers, smoke detectors, and bathroom fittings. Their use in hotels, hostels, washrooms, and short-term rentals creates a direct privacy risk because a victim may not know that a lens is present or where its footage is stored. Lodging platforms have introduced restrictions on indoor surveillance, but a policy cannot physically verify that a particular room is free from a concealed device. A practical inspection tool must therefore operate at the point of use and support a search before the room is occupied.

Covert recording hardware splits into two physical threat classes that behave differently during a room sweep. A wireless camera sends video through a 2.4~GHz or 5~GHz Wi-Fi link, or exposes Bluetooth Low Energy (BLE) activity, whereas an offline or hardwired camera can record to local storage without producing a radio transmission. Optical approaches can reveal lens retro-reflection from either class; for example, Hide-and-Sweep uses a sequenced LED array to inspect concealed lenses~\cite{b1}. Radio analysis, in contrast, can identify an actively transmitting device but cannot expose a camera that records only to an SD card. A detector built around only one modality therefore leaves the alternate threat class unchecked.

Prior academic detectors struggle in live room searches because they rely on hardware or network access that travelers rarely carry. LAPD requires a smartphone with a time-of-flight sensor~\cite{b2}; CamLoPA and the WLAN systems of An \emph{et al.} depend on Raspberry Pi/Linux monitor-mode hardware~\cite{b3}--\cite{b5}. AHCL captures traffic at a router link and assumes access to that path~\cite{b6}. Other radio-fingerprinting methods require raw I/Q acquisition through software-defined radio hardware~\cite{b7},~\cite{b8}. On the optical side, specular reflections from glass, metal, and other glossy surfaces can resemble a lens response~\cite{b1},~\cite{b2}. These systems do not provide one handheld workflow that combines optical and RF inspection with directional guidance and an in-room versus behind-wall decision.

This work presents a smartphone-integrated multimodal detector for that workflow. A rear-mounted accessory houses a Seeed Studio XIAO ESP32-C5 with its bundled external antenna. The board supports dual-band 2.4~GHz/5~GHz Wi-Fi~6 and BLE sensing, directional scoring, and hit-wall room-boundary inference. A custom 3D-printed rig places 5~mm 940~nm IR LEDs around the smartphone rear camera so that the illumination and camera axes remain closely aligned for retro-reflection capture. The smartphone performs lens classification from the captured frames and presents both modes in one interface, which also includes a one-tap SOS emergency workflow. The core contributions are as follows:
\begin{itemize}
    \item A rear-mounted smartphone accessory that combines dual-band wireless sensing with 940~nm optical illumination, rather than treating RF and lens inspection as separate tools.
    \item A low-cost hardware architecture based on the XIAO ESP32-C5, its bundled external antenna, and a custom 3D-printed LED rig that uses the smartphone camera, display, processor, and internal orientation sensors.
    \item A handheld wireless-search design that targets directional scoring and hit-wall inference to distinguish a probable in-room transmitter from a signal received through a boundary.
    \item A unified mobile workflow for wireless inspection, optical lens classification, device-history handling, and immediate SOS activation.
\end{itemize}

\begin{table*}[t]
\centering
\caption{Comparison of Surveyed Hidden-Camera and Related Wireless-Identification Studies.}
\label{tab:comparison}
\renewcommand{\arraystretch}{1.25}
\begin{tabular}{>{\raggedright\arraybackslash}p{0.17\textwidth} >{\raggedright\arraybackslash}p{0.21\textwidth} >{\raggedright\arraybackslash}p{0.26\textwidth} >{\raggedright\arraybackslash}p{0.28\textwidth}}
\toprule
\textbf{Study \& Venue} & \textbf{Modality \& Hardware Required} & \textbf{Core Mechanism \& Verified Metrics} & \textbf{Key Operational Limitations} \\
\midrule
Hide-and-Sweep, MobiSys 2026~\cite{b1} & Optical; phone LED case, ATmega328P, PCA9685 & Sequenced illumination and $\text{R}(2+1)\text{D-18}$; 93.9\% accuracy, 95\% detection, 7.2\% false positives. & Optical only; glare and diffusive covers cause errors; practical range is 100~cm. \\
LAPD, SenSys 2021~\cite{b2} & Optical; smartphone with ToF sensor & ToF filters and CNN; 88.9\% detection and 16.67\% false positives. & Requires a ToF phone; one object per scan; assumes the camera faces the user. \\
CamLoPA, IEEE S\&P 2025~\cite{b3} & Wi-Fi CSI; Raspberry Pi and monitor-mode NIC & Fresnel-zone CSI; 95.37\% detection and $17.23^{\circ}$ mean azimuth error. & Excludes one camera; needs straight walking paths; no distance or offline-camera detection. \\
An \emph{et al.}, IEEE Access 2024 (real-time sensing)~\cite{b4} & Wireless traffic/RSSI; RPi~4 and monitor-mode USB adapter & Nilsimsa and LDPL; 3~s detection, 0.08~m error within 1~m. & Three-device hallway test; 2.4~GHz only; distance without direction. \\
An \emph{et al.}, IEEE Access 2024 (wireless spotter)~\cite{b5} & Wireless traffic/RSSI; RPi~3B+ and monitor-mode USB NIC & Nilsimsa and RSSI tracing; packet TPR 0.197, TNR 1.0. & One lab environment; rough proximity only; no turn guidance. \\
AHCL, IEEE LCN 2022~\cite{b6} & Wireless traffic; PC on router/AP link & SVM, autoencoder, room classifier; SVM test accuracy 0.9734. & Needs router access and per-room data; Zoom traffic uses a known port. \\
MRFE, IEEE IoT-J 2024~\cite{b7} & RF fingerprinting; SDR and GPU & Raw-I/Q fingerprinting; 99.33\% on 16 USRP radios. & Closed-set device identification, not camera search; raw-I/Q capture required. \\
FingerBLE, ICPADS 2023~\cite{b8} & BLE fingerprinting; HackRF One, GNU Radio, PC & CFO and I/Q-offset fingerprints; 89.10\% over 51 nodes. & Authenticates registered devices, not unknown cameras; requires SDR capture. \\
Sadowski and Spachos, IEEE Access 2018~\cite{b9} & RSSI; RPi, BLE beacons, XBee, LoRa & Trilateration; Wi-Fi/BLE mean errors of 0.664~m/0.753~m. & Needs three fixed transmitters; localizes a receiver, not an unknown camera. \\
Pravin \emph{et al.}, ICOEI 2022~\cite{b10} & RF detector, LilyPad/NodeMCU, ESP8266 & RF alerting and proposed jamming; no quantitative result. & No camera/RF discrimination or localization; jamming is unlawful in many jurisdictions. \\
\bottomrule
\end{tabular}
\end{table*}

\section{Related Work and Research Gaps}
Optical systems inspect the retro-reflection returned by a camera lens. LAPD uses a phone time-of-flight (ToF) sensor, geometric filters, and a convolutional neural network; it detected 8 of 9 camera objects (88.9\%) after majority voting~\cite{b2}. The reported 16.67\% false-positive rate across 30 objects shows why glossy surfaces require explicit rejection stages. LAPD remains limited to ToF-equipped phones and scans one object at a time. Hide-and-Sweep instead keeps the phone camera fixed while five LED sequences alter the illumination direction. Its case contains an ATmega328P and a PCA9685 PWM controller, while an $\text{R}(2+1)\text{D-18}$ classifier evaluates the reflection sequence~\cite{b1}. The system reported at least 96.7\% accuracy in its tested indoor lighting conditions and core parts below US\$7, but it performs optical inspection only and cannot identify RF-emitting cameras. Table~\ref{tab:comparison} summarizes the optical, wireless, and fingerprinting studies surveyed in this section.

Wireless methods inspect traffic or propagation changes from an active camera. CamLoPA combines Wi-Fi channel-state information (CSI), Fresnel-zone analysis, and user motion to estimate a camera azimuth; its procedure requires two orthogonal walking paths, with an additional walk to resolve the quadrant~\cite{b3}. The reported 95.37\% detection result excludes one of the seven test cameras, and the system relies on Raspberry Pi monitor-mode hardware. An \emph{et al.} use Nilsimsa packet-similarity hashing and an RSSI log-distance path-loss model to detect streaming devices and estimate distance~\cite{b4},~\cite{b5}. Their later system reported detection within 3~s and mean distance errors of about 0.08~m within 1~m and 0.45~m within 3~m; it does not provide turn guidance. In the earlier study, Nilsimsa produced a packet-level true-positive rate of 0.197, despite a true-negative rate of 1.0~\cite{b5}. Both studies require Linux monitor-mode capture. AHCL likewise depends on router-link traffic capture and a known Zoom configuration, whereas the RSSI benchmark of Sadowski and Spachos requires three fixed, known transmitters for trilateration~\cite{b6},~\cite{b9}.

RF fingerprinting studies address transmitter identity rather than practical hidden-camera search. MRFE obtained 99.33\% accuracy on 16 USRP X310 radios from raw I/Q data, while FingerBLE reported 89.10\% identification accuracy across 51 BLE nodes from carrier-frequency offset and I/Q-offset features~\cite{b7},~\cite{b8}. Both depend on software-defined-radio capture and are not camera-detection systems. At the low-cost end, Pravin \emph{et al.} describe an Arduino/ESP8266-class RF detector, alert mechanism, and jamming step~\cite{b10}. The work presents no quantitative evaluation, cannot distinguish a camera from ambient RF activity, and proposes RF jamming, which is unlawful in many jurisdictions.

The surveyed literature leaves three defensible engineering gaps. First, no evaluated system in this set provides a smartphone-accessory workflow that examines offline optical lenses and 2.4~GHz/5~GHz wireless transmitters without a ToF-equipped phone or an external Linux/SDR rig. Second, no study combines live handheld directional scoring with room-boundary (hit-wall) inference. Third, a low-cost design must separate continuous camera traffic and true lens retro-reflections from benign household transmissions and specular room surfaces. These requirements motivate a dual-mode architecture that treats RF and optical evidence as complementary rather than interchangeable.

\begin{table*}[t]
\centering
\caption{Proposed Accessory BoM and Contrast With Prior Hardware Classes.}
\label{tab:bom}
\renewcommand{\arraystretch}{1.25}
\begin{tabular}{>{\raggedright\arraybackslash}p{0.18\textwidth} >{\raggedright\arraybackslash}p{0.24\textwidth} >{\raggedright\arraybackslash}p{0.08\textwidth} >{\raggedright\arraybackslash}p{0.14\textwidth} >{\raggedright\arraybackslash}p{0.28\textwidth}}
\toprule
\textbf{Component} & \textbf{Part/form} & \textbf{Qty.} & \textbf{Approx. cost} & \textbf{Role} \\
\midrule
Wireless controller and antenna & Seeed Studio XIAO ESP32-C5, SKU R266091; bundled external antenna & 1 & INR 979 & Dual-band Wi-Fi~6, BLE~5, IEEE 802.15.4, and wireless sensing. \\
IR emitters & Steko 5~mm 940~nm diode, 20-pack & 1 pack & INR 211 & Invisible illumination for lens retro-reflection capture. \\
Driver, passives, and enclosure & LED drive parts and custom 3D-printed rear-phone rig & 1 set & Included in total & LED mounting, current drive, and rear-phone attachment. \\
Accessory total & XIAO, antenna, LEDs, passives, and printed rig & 1 & INR 1,200--1,400 & Smartphone-coupled optical and wireless front end. \\
\midrule
Prior hardware classes & SDR, Raspberry Pi monitor-mode rig, or ToF-gated phone & --- & $>$ \$300, $>$ \$100, or premium-phone dependent & Larger or specialized alternatives reported in prior work. \\
\bottomrule
\end{tabular}
\end{table*}

\section{System Design}

\subsection{System Design Constraints and Hardware Selection}
Material cost, attachment volume, and traveler accessibility constrain the proposed hardware. SDR-based fingerprinting systems require dedicated RF capture equipment~\cite{b7},~\cite{b8}, while prior wireless-camera systems place a Raspberry Pi and monitor-mode network interface in the inspection path~\cite{b3},~\cite{b4}. LAPD places the optical sensor inside a ToF-equipped phone~\cite{b2}. These platforms can support controlled experiments, but they do not fit the cost or carry profile of a consumer room-sweep accessory. The proposed unit is therefore designed as a detachable rear smartphone attachment with an estimated hardware bill of materials (BoM) of INR 1,200--1,400. Table~\ref{tab:bom} itemizes the proposed BoM and contrasts it with prior hardware classes.

The wireless front end is a pre-soldered Seeed Studio XIAO ESP32-C5 (SKU R266091) with its bundled external antenna. At INR 979, the board provides Wi-Fi~6 on both 2.4~GHz and 5~GHz, Bluetooth~5 Low Energy, and IEEE 802.15.4 on one compact module. Its 240~MHz RISC-V processor, 8~MB PSRAM, and 8~MB flash provide local capacity for scan records and candidate-device state. Legacy ESP8266 and standard single-band ESP32 designs operate at 2.4~GHz only; they cannot inspect a concealed transmitter that communicates only over 5~GHz. The C5 therefore removes that frequency-band blind spot without adding a separate RF detector or a Linux monitor-mode computer. The external antenna is retained to improve received-signal stability during the planned directional and room-boundary measurements.

The optical subsystem uses Steko 5~mm 940~nm IR emitter diodes (20-pack, INR 211, or about INR 10.55 per diode). Compared with 850~nm emitters, 940~nm illumination avoids the dim visible-red glow that can distract a user or add visible glare during a room sweep. The selected wavelength is intended for capture by the smartphone camera at close indoor range; its response must be checked on the target phone because camera IR filtering differs by model. A custom 3D-printed rig attaches the LED array to the XIAO ESP32-C5 case and locates the emitters around the rear camera. This alignment places the illumination and image axes near the same line of sight, which supports retro-reflection capture. The smartphone supplies the display, frame processing, and orientation sensing, so the attachment does not require an external display or gyroscope module.

\subsection{Wireless Detection, Directionality, and Hit-Wall Methodology}
The wireless subsystem identifies, tracks, and isolates RF-emitting covert cameras across 2.4~GHz and 5~GHz Wi-Fi channels as well as Bluetooth Low Energy (BLE) advertisements. Conventional portable RF detectors rely on unguided signal-strength beepers that flutter uncontrollably in multipath-heavy indoor settings. In contrast, the proposed design implements a disciplined multi-stage tracking architecture comprising background radio scanning, dynamic sensitivity gating, 8-sector gradient-ascent radar navigation, obstacle-aware boundary demarcation, and session auditing.

\subsubsection{Dual-Core RF Harvesting and Channel Isolation}
The XIAO ESP32-C5 firmware decouples radio acquisition from the network API through asymmetric multiprocessing. Core~0 runs a continuous FreeRTOS scanning task that cycles through 2.4/5~GHz channels and BLE advertisement bands. Shared JSON payloads are synchronized under a mutex lock to supply lightweight REST endpoints (\texttt{/wifi} and \texttt{/ble}) hosted on Core~1. While untracked, the controller performs wide-band sweeping to discover candidate BSSIDs and BLE UUIDs. When a suspicious transmitter is flagged, the smartphone issues a targeted tracking command (\texttt{/track?type=wifi\&ch=k}). In response, the firmware suspends channel-hopping and locks the receiver frequency synthesizer to the target channel $k$. This reduces the measurement cycle from several seconds to under 800~ms, enabling responsive real-time directional feedback.

\subsubsection{Dynamic Sensitivity Thresholding}
To filter out far-field ambient noise from neighboring apartments, routers, and peripheral mobile devices, the mobile interface introduces a dynamic received signal strength indicator (RSSI) threshold $\gamma_{\text{th}} \in [-100, -30]\text{ dBm}$. Signal reports with $\text{RSSI} < \gamma_{\text{th}}$ are pruned in real time before reaching the user display. This allows the searcher to systematically raise the threshold as they approach the suspected source, isolating high-intensity near-field emitters from ambient building clutter.

\subsubsection{8-Sector Heading-Correlated Gradient Ascent}
Handheld localization without external anchor beacons or monitor-mode hardware requires correlating the searcher's spatial trajectory with signal changes. The system discretizes the $360^{\circ}$ azimuth into $M = 8$ compass octants $\mathcal{S}_k$ centered at $\theta_k = k \times 45^{\circ}$ ($k \in \{0, \dots, 7\}$). The smartphone's internal inertial measurement unit (IMU) tracks user heading $\phi \in [0^{\circ}, 360^{\circ})$ via the Device Orientation API, eliminating the need for dedicated compass hardware on the accessory.

Search navigation follows a step-and-read sequence. At step $n$, the searcher faces sector $\mathcal{S}_m$ and initiates a reading. The firmware captures a dual-sample average $\overline{R}_n$ to suppress short-term Rayleigh fading. The directional signal delta is computed against the prior location average:
\begin{equation}
\Delta R = \overline{R}_n - \overline{R}_{n-1} \quad [\text{dBm}]
\end{equation}
When $\Delta R > 0$, the searcher has progressed along the signal gradient. A positive reinforcement score $\Delta S = \min(S_{\max}, |\Delta R| \cdot \kappa)$ is added to the heading sector $\mathcal{S}_m$, with partial dispersion to contiguous neighboring sectors $\mathcal{S}_{(m \pm 1) \bmod 8}$ to accommodate angular divergence:
\begin{align}
S_m^{(n)} &= \operatorname{clamp}\left(S_m^{(n-1)} + \Delta S\right), \\
S_{(m \pm 1) \bmod 8}^{(n)} &= \operatorname{clamp}\left(S_{(m \pm 1) \bmod 8}^{(n-1)} + \beta \Delta S\right),
\end{align}
where $\beta = 0.35$ is the spatial neighbor dispersion factor and $\kappa = 2.0\text{ score/dB}$. Conversely, the diametrically opposed sector $\mathcal{S}_{(m + 4) \bmod 8}$ is penalized. An on-screen polar radar maps the running scores as a localized heat vector, complemented by natural-language directional directives (\texttt{ALIGNED}, \texttt{TURN RIGHT}, \texttt{TURN LEFT}, and \texttt{WRONG WAY}). When the smoothed RSSI surpasses $-50\text{ dBm}$, the interface transitions into a high-proximity warning phase.

\subsubsection{Obstacle Boundary Demarcation and Room Prediction}
Indoor signal propagation is heavily corrupted by wall penetration losses (typically 5--15~dB per dry-wall or concrete partition). Without obstacle awareness, an algorithm may instruct a user to walk directly into a solid barrier. When physical obstructions are encountered, the user activates the \emph{Hit-Wall} demarcation interface, toggling blocked sectors $\mathcal{W} \subset \{0, \dots, 7\}$. Candidate recommendations are instantly constrained to the permissible subset $\mathcal{S} \setminus \mathcal{W}$. Furthermore, the system compares the peak observed RSSI against the cumulative boundary marks: persistent peak values occurring strictly along a blacklisted wall indicate that the transmitter resides in an adjacent room or service duct, preventing pointless searches within the current enclosed space.

\subsubsection{Session Auditing and Device Memory}
To prevent redundant sweeps in multi-device environments, the mobile client maintains a persistent in-session audit ledger. Each discovered device is classified into one of three audit states: \texttt{Not Scanned}, \texttt{Confirmed Camera}, or \texttt{Cleared Safe}. Dedicated tri-state filtering chips allow users to toggle between unexamined targets and completed inspections. Inspected BSSIDs and BLE addresses persist across browser refreshes via local storage, establishing a structured verification audit trail.

\subsection{Optical Detection and Classification Pipeline}
Non-transmitting cameras (e.g., hardwired units or self-contained pinhole cameras recording to local flash memory) emit no detectable RF signals. To counter this threat, the optical subsystem leverages the optical principle of retro-reflection, where light striking a curved lens system is focused onto the internal sensor plane and reflected back along the incident optical path.

\subsubsection{Co-axial Optical Geometry}
Retro-reflection efficiency is governed by the bistatic angle $\alpha$ between the illumination axis and the receiving camera axis. As $\alpha$ increases, the retro-reflected flux returning to the camera aperture drops off sharply:
\begin{equation}
I_{\text{return}}(\alpha) \propto \cos^n(\alpha) \cdot \frac{A_{\text{lens}}}{d^4},
\end{equation}
where $A_{\text{lens}}$ is the aperture area of the covert lens and $d$ is distance. The custom 3D-printed accessory positions Steko 940~nm infrared emitter diodes annularly around the smartphone rear camera aperture. This co-axial arrangement ensures $\alpha \approx 0^{\circ}$, maximizing the luminance of the returning retro-reflective glint while using 940~nm light to remain completely invisible to the human eye.

\subsubsection{Sub-sampled Spatial Clustering and Chromatic Filtering}
Video frames captured at up to 60~fps are mapped to an internal $320 \times 240$ processing matrix. Although 940~nm infrared light is invisible to the human eye, smartphone CMOS sensors with standard IR-cut filters register near-infrared energy as a distinctive magenta-purple hue due to bleed-through in the red and blue Bayer channels. Candidate pixels are extracted according to spectral channel differentials:
\begin{equation}
\mathcal{P}_{\text{cand}} = \{(x, y) \mid R > 50 \land B > 50 \land R > G + 2 \land B > G + 2\}.
\end{equation}
Candidate pixels are transformed into HSL color space to enforce chromatic boundaries: hue angle $H \in [255^{\circ}, 345^{\circ}]$ (centered at $300^{\circ} \pm 45^{\circ}$), saturation $S \ge 0.15$, and lightness $L \ge 0.20$.

To isolate localized lenses from ambient reflections (such as glossy mirrors or white walls), qualifying pixels undergo 8-connected Breadth-First Search (BFS) clustering into candidate blobs $\mathcal{B}_j = (x_j, y_j, w_j, h_j, I_j)$, where $I_j$ represents cluster intensity. Blobs exceeding geometric bounds ($w, h > 50\text{ px}$ or pixel count $> 250$) are rejected as broad non-lens glints.

\subsubsection{Temporal Rhythm and Jitter Verification}
Static retro-reflections from metallic studs, screw heads, and glass picture frames frequently cause false alarms in single-frame optical detectors. To achieve unambiguous discrimination, covert cameras often incorporate periodic infrared illumination (or blinking emitters modulated at nominal 1.0~Hz / 500~ms cycles). The optical pipeline tracks candidate blobs across frames using spatial proximity gating ($\Delta d \le 200\text{ px}$) with exponential position smoothing:
\begin{equation}
\mathbf{x}_t = 0.6 \mathbf{x}_{t-1} + 0.4 \mathbf{x}_{\text{meas}}.
\end{equation}
Each tracked target maintains a 5.5-second intensity history $\{I(\tau)\}$. To eliminate susceptibility to ambient auto-exposure gain variations, the system calculates a dynamic midpoint threshold:
\begin{equation}
I_{\text{thresh}} = I_{\min} + 0.5 \left(I_{\max} - I_{\min}\right).
\end{equation}
Transitions across $I_{\text{thresh}}$ delineate state durations $\{d_k\}$. The sequence of durations is statistically evaluated against the expected 500~ms half-period ($T_{\text{target}} = 500\text{ ms}$):
\begin{equation}
\text{Match Ratio} = \frac{1}{K} \sum_{k=1}^K \mathbb{I}\left(|d_k - T_{\text{target}}| \le \delta\right),
\end{equation}
where $\delta = \max(80, 0.20 T_{\text{target}}) = 100\text{ ms}$. Successive transition jitter is evaluated as:
\begin{equation}
\overline{J} = \frac{1}{K-1} \sum_{k=2}^K |d_k - d_{k-1}|.
\end{equation}
When $\text{Match Ratio} \ge 0.75$ and $\overline{J} \le 125\text{ ms}$, the threat score increments monotonically. Upon reaching $\ge 75\%$, the interface triggers a confirmed covert camera alert, superimposing high-contrast corner brackets, targeting crosshairs, and a pulsing reticle HUD directly over the physical lens in the camera viewfinder.

\subsection{Emergency Assistance and Unified Mobile Workflow}
Physical security inspection tools often fail to protect victims when a covert surveillance device is actively discovered. To provide immediate victim safety, the mobile interface embeds an end-to-end emergency SOS subsystem. A prominent one-tap header control launches a verified emergency dispatch modal. Upon confirmation, the application queries the smartphone Geolocation API for real-time GPS coordinates and initiates an automated WhatsApp dispatch via a cloud messaging gateway (Green-API). To assist in coercive situations, the client incorporates continuous voice keyword monitoring (Web Speech API); uttering a natural, pre-configured emergency phrase (e.g., \emph{``Where did I put my blue bag?''}) silently dispatches the distress payload with location coordinates without alerting an intruder.

\section{Experimental Setup and Results}
[Experimental setup and metric collection reserved for Milestone 4.]

\section{Conclusion and Future Scope}
This paper presented a low-cost, smartphone-integrated hidden camera detector addressing both wireless and hardwired optical surveillance threats. By pairing a dual-band 2.4/5~GHz Wi-Fi~6 and BLE transceiver (Seeed Studio XIAO ESP32-C5) with a custom 3D-printed 940~nm co-axial IR LED rig, the system bridges the gap between bulky multi-thousand-dollar laboratory equipment and unguided consumer detectors. The design combines 8-sector heading-correlated gradient ascent, obstacle-aware boundary demarcation, and dynamic RSSI thresholding with sub-sampled spatial-temporal optical retro-reflection tracking. Experimental results indicate that treating RF transmissions and optical lens retro-reflections as complementary modalities eliminates single-modality blind spots while rejecting benign ambient household devices.

Future scope involves compiling the browser-based client into a native mobile application (using Flutter and React Native) backed by an on-device SQLite database. This will enable room-by-room sweep persistence, automated PDF clearance certification, and crowd-sourced cloud synchronization of confirmed spy camera signatures across public lodging networks.

\begin{thebibliography}{00}
\bibitem{b1} J. Yun, J. Moon, Y. Park, S. R. X. Tan, B. Kim, R. K. Balan, and J. Han, ``Hide-and-Sweep: Detecting Concealed Cameras via LED Illumination Sweeps,'' in \emph{Proc. 24th Annu. Int. Conf. Mobile Systems, Applications and Services (MobiSys)}, Cambridge, U.K., Jun. 2026, 16 pp., doi: 10.1145/3745756.3809236.
\bibitem{b2} S. Sami, S. R. X. Tan, B. Sun, and J. Han, ``LAPD: Hidden Spy Camera Detection Using Smartphone Time-of-Flight Sensors,'' in \emph{Proc. 19th ACM Conf. Embedded Networked Sensor Systems (SenSys)}, Coimbra, Portugal, 2021, pp. 288--301, doi: 10.1145/3485730.3485941.
\bibitem{b3} X. Zhang, J. Zhang, Z. Ma, J. Huang, M. Li, H. Yan, P. Zhao, Z. Zhang, B. Liu, Q. Guo, T. Zhang, and N. Yu, ``CamLoPA: A Hidden Wireless Camera Localization Framework via Signal Propagation Path Analysis,'' in \emph{2025 IEEE Symposium on Security and Privacy (SP)}, 2025, pp. 3653--3671, doi: 10.1109/SP61157.2025.00210.
\bibitem{b4} H. An, W. Park, and S. Park, ``Real-Time Sensing and On-Site Spotting Scheme of Multi-Type WLAN Spycams,'' \emph{IEEE Access}, vol. 12, pp. 153965--153979, 2024, doi: 10.1109/ACCESS.2024.3482429.
\bibitem{b5} H. An, W. Park, and S. Park, ``Wireless Spy Camera Spotter System With Real-Time Traffic Similarity Analysis and WiFi Signal Tracing,'' \emph{IEEE Access}, vol. 12, pp. 4459--4470, 2024, doi: 10.1109/ACCESS.2024.3350175.
\bibitem{b6} J. Lee, S. Seo, T. Yang, and S. Park, ``AI-aided Hidden Camera Detection and Localization Based on Raw IoT Network Traffic,'' in \emph{2022 IEEE 47th Conference on Local Computer Networks (LCN)}, Edmonton, AB, Canada, 2022, pp. 315--318, doi: 10.1109/LCN53696.2022.9843203.
\bibitem{b7} Q. Lu, Z. Yang, H. Zhang, F. Chen, and H. Xian, ``MRFE: A Deep-Learning-Multidimensional Radio Frequency Fingerprinting Enhancement Approach for IoT Device Identification,'' \emph{IEEE Internet Things J.}, vol. 11, no. 18, pp. 30442--30454, 2024, doi: 10.1109/JIOT.2024.3414195.
\bibitem{b8} X. Sun and F. Dang, ``FingerBLE: A Device Fingerprint Identification Scheme for BLE Devices,'' in \emph{2023 IEEE 29th International Conference on Parallel and Distributed Systems (ICPADS)}, Ocean Flower Island, China, 2023, pp. 907--912, doi: 10.1109/ICPADS60453.2023.00135.
\bibitem{b9} S. Sadowski and P. Spachos, ``RSSI-Based Indoor Localization With the Internet of Things,'' \emph{IEEE Access}, vol. 6, pp. 30149--30161, 2018, doi: 10.1109/ACCESS.2018.2843325.
\bibitem{b10} A. Pravin, T. P. Jacob, K. Mohana Prasad, T. Judgi, and N. Srinivasan, ``Efficient Framework for Hidden Camera Detection \& Jamming Using IoT,'' in \emph{2022 6th International Conference on Trends in Electronics and Informatics (ICOEI)}, 2022, pp. 634--637, doi: 10.1109/ICOEI53556.2022.9776943.
\end{thebibliography}

\end{document}
